%%%
\chapter{User Model Study} \label{proposed_study}
%%%

\section{Purpose}

Accurate evaluation of the various synchronization algorithms within the simulation environment is dependant upon a set of accurate user models.  These models dictate the frequency of record modification, data synchronization, and the users daily movements throughout their environment.  An accurate model will thus provide a more useful result for each algorithm, and thus a more accurate cost comparison can be made.

Building accurate user models requires information on the habits and patterns of users with handheld mobile devices which are synchronized against a database.  Unfortunately, such habits and patterns are lacking within the existing literature.  In addition, due to the fast evolution experienced in the mobile device arena, any possible existing studies older than one or two years may not reflect the habits of today's mobile device user.  As such, a study of mobile device user habits was required in order to build a series of accurate, real-world user models to evaluate the simulated synchronization algorithms.

\section{Methodology}

In order to study the habits of the target population for the purpose of creating accurate user models, a series of seven survey questions was derived.  These questions are listed in Appendix \ref{modelquestions}.  The questions are designed to determine the average number of times users are expected to synchronize during both work hours and off-work hours, along with the size of the database being synchronized against and the average number of changes to said database during an average day.  Four of the questions are designed to elicit specific numerical responses from which to design a user model, while the remaining three questions are designed to determine the respondent's confidence level for the values reported in three of the other four questions.

The survey itself is presented as a web application, using LimeSurvey \cite{limesurvey}.  This online survey tool was installed on a secure server at the University of Victoria, and was used to manage the survey invitations, provide the facility for the target population to complete the survey via the World-Wide Web, and to anonymize and store the survey results.  Prior to completing the survey, the users are invited to download and read a \underline{Letter of Implied Consent}, which is available in Appendix \ref{impliedconsent}.

Participants were self-selected, with invitations being sent to organizations known or believed to use data synchronization software.  This was achieved by targeting Internet mailing lists whose purpose is the discussion of implementing mobile data synchronization solutions in multi-user organizations.  Organizations known to use the jSyncManager \cite{RefWorks:30} were also invited to participate.  A copy of the invitation letter can be found in Appendix \ref{surveyinvitation}.  In order to minimize responses from organizations or individuals outside the target population, each invitation contained a survey-software generated token, which must be supplied in order to respond to the survey.  Each token may only be used once, and no information is stored or retained to attach a specific invitee to a specific token.

The study was reviewed and approved by the University of Victoria Human Research Ethics Board.  A copy of the Certificate of Approval can be found in Appendix \ref{ethicscert}.

\section{Results}

Raw results are listed in Appendix \ref{surveyresults}.

\section{Conclusion}

Concluding remarks.
